% gerado por gera_dados.py (la_mef): nao editar a mao
\begin{table}[H]\centering\footnotesize\setlength{\tabcolsep}{4pt}
\caption{Os modos de ruptura, em \% (como na Fig.~20 do Ceron): pelo pé, pela face (acima do pé) e pela base (abaixo do pé). LA: o mecanismo crítico de cada realização (sem peso na ISMC; no direto, sobram as com nenhum mecanismo admissível, até 0,2\,\%). MEF: pela banda de cisalhamento (verificações de 21/11, V07; a banda sempre pega o pé, então uma saída a poucos centímetros acima conta como pé).}\label{tab:lamodos}
\begin{tabular}{llrrrrrrr}\toprule
 & & & \multicolumn{3}{c}{todas as realizações} & \multicolumn{3}{c}{as que falham} \\\cmidrule(lr){4-6}\cmidrule(lr){7-9}
$r$ & caso & $N$ (falhas) & pé & face & base & pé & face & base \\\midrule
$0$ & LA, MC direto & 10\,460 (409) & 74,8 & 24,8 & 0,3 & 74,1 & 25,9 & 0,0 \\
 & Ceron (2025), Fig.~20 (LA) & 50\,000 & 80,8 & 19,2 & 0,0 & -- & -- & -- \\
$-0{,}5$ & LA, MC direto & 60\,250 (403) & 83,9 & 15,7 & 0,2 & 87,1 & 12,4 & 0,5 \\
 & LA, ISMC s21 (Nyström, 4 centros) & 10\,000 & -- & -- & -- & 90,8 & 8,9 & 0,3 \\
 & MEF, MC direto (V07) & 24\,000 (408) & 98,0 & 0,1 & 1,9 & 99,5 & 0,0 & 0,5 \\
$-0{,}7$ & LA, MC direto (parcial) & 7\,930 (9) & 88,5 & 11,1 & 0,3 & 100,0 & 0,0 & 0,0 \\
 & LA, ISMC s21 (Nyström, 4 centros) & 10\,000 & -- & -- & -- & 95,3 & 4,5 & 0,2 \\
 & MEF, MC direto (V07) & 110\,000 (403) & 99,3 & 0,0 & 0,7 & 99,8 & 0,0 & 0,2 \\
\bottomrule\end{tabular}\end{table}
